\documentclass[12pt]{article}
\usepackage{fancyhdr,listings,amsfonts,amsmath,color,lastpage}
\usepackage{float}
\usepackage{multirow} 
\usepackage{enumitem}
\usepackage{caption}
\usepackage{tikz}
\usetikzlibrary{arrows.meta, positioning, shapes.geometric}

\tikzset{
    node style/.style={draw, rectangle, rounded corners, minimum width=3.5cm, minimum height=1.2cm, text centered, font=\sffamily},
    decision style/.style={draw, diamond, aspect=2, minimum width=3.5cm, text centered, font=\sffamily},
    arrow style/.style={->, thick},
}

% code listing settings
\usepackage{listings}
\lstset{
    language=Python,
    basicstyle=\ttfamily\small,
    aboveskip={1.0\baselineskip},
    belowskip={1.0\baselineskip},
    columns=fixed,
    extendedchars=true,
    breaklines=true,
    tabsize=4,
    prebreak=\raisebox{0ex}[0ex][0ex]{\ensuremath{\hookleftarrow}},
    frame=lines,
    showtabs=false,
    showspaces=false,
    showstringspaces=false,
    keywordstyle=\color[rgb]{0.627,0.126,0.941},
    commentstyle=\color[rgb]{0.133,0.545,0.133},
    stringstyle=\color[rgb]{01,0,0},
    numbers=left,
    numberstyle=\small,
    stepnumber=1,
    numbersep=10pt,
    captionpos=t,
    escapeinside={\%*}{*)}
}

\pagestyle{fancy}
\parindent0em

\newcommand{\masunitnumber}{CSC111}
\newcommand{\examdate}{ JULY 2026}
\newcommand{\academicyear}{2025/2026}
\newcommand{\semester}{I}
\newcommand{\university}{UNIVERSITY OF ESWATINI}
\newcommand{\faculty}{FACULTY OF SCIENCE AND ENGINEERING}
\newcommand{\department}{DEPARTMENT OF COMPUTER SCIENCE}

\newcommand{\coursename}{Introduction to Computer Science} 
\newcommand{\numberofhours}{3}
\newcommand{\Z}{\mathbb{Z}}
\newcommand{\F}{\mathbb{F}}
\newcommand{\C}{\mathbb{C}}
\newcommand{\Q}{\mathbb{Q}}
\newcommand{\PP}{\mathbb{P}}
\newcommand{\R}{\mathbb{R}}
\newcommand{\GL}{\mathrm{GL}}
\definecolor{listinggray}{gray}{0.9} 
\lstset{backgroundcolor=\color{listinggray},rulecolor=\color{blue}}
\lstset{linewidth=\textwidth} 
\lstset{commentstyle=\tt, stringstyle=\ttfamily,showspaces=false,basicstyle=\tt}
\lstset{showstringspaces=false, keepspaces=true}
\newcommand{\code}[1]{\lstinline@#1@}
\newcommand{\Code}[1]{\colorbox{listinggray}{\lstinline@#1@}}

\lhead{}
\rhead{}
\chead{}  % Removed university name from header
\lfoot{}
\rfoot{}
\cfoot{\thepage\ of \pageref{LastPage}}

\begin{document}
\setlength{\headsep}{14truemm}
\setlength{\headheight}{14truemm}
\setlength{\voffset}{-0.4truein}
\renewcommand{\headrulewidth}{0.0pt}

% First page with university name in title block
 



\begin{center}
 
 \textbf{\large \university }  \\
  \textbf{\large \faculty } 
    \textbf{\large \department }  

\bf SEMESTER \semester\ RESIT EXAMINATION \\ \academicyear
\end{center}
\begin{center}
{\bf \masunitnumber -- \coursename}
\end{center}
\vspace{20truemm}

\noindent \examdate\hspace{55truemm} TIME ALLOWED: \numberofhours\ HOURS

\vspace{10truemm}
\hrule

\begin{minipage}{\textwidth}
\vspace{2} 
\small{\textbf{Target Group(s):}} \hfill   \textbf{Marks=100} \\
\textbf{B.Sc. Information Technology (FT/IDE)}, \textbf{B.Sc. Computer Science Education (FT/IDE)}, \textbf{B.Eng.}, \textbf{B.Sc. Information Science}, \textbf{B.Sc. Geographic Information Science}
\end{minipage}

\vspace{10truemm}
\noindent\underline{INSTRUCTIONS}
\vspace{5mm}

\begin{enumerate}
    \item This examination paper contains {\bf THREE (3)} sections and comprises {\bf SEVEN (7)} printed pages including the cover page.
    \item Answer \textbf{ALL} questions in \textbf{Section A} (50 marks).
    \item Answer \textbf{EITHER} all questions in \textbf{Section B} (50 marks) \textbf{OR} all questions in \textbf{Section C} (50 marks).
    \item Answer each section beginning on a {\bf FRESH} page of the answer book.
    \item Any type of calculator is not allowed.
\end{enumerate}

\newpage

% Subsequent pages with no university name in header
\section*{SECTION A: COMPULSORY (50 marks)}
\textit{Answer all questions in this section.}

\begin{enumerate}[left=0pt, label=\bfseries Q\arabic*.]
    \item \textbf{Computer Science Disciplines} \hfill (10 marks)
    \begin{enumerate}[label=(\alph*)]
        \item Define Computer Science and explain why it is considered more than just programming. \hfill (4 marks)
        \item List and briefly describe three other distinct computing disciplines identified by the Association for Computing Machinery (ACM). \hfill (6 marks)
    \end{enumerate}

    \item \textbf{Computer Generations} \hfill (10 marks)
    \begin{enumerate}[label=(\alph*)]
        \item Describe the technological breakthrough that characterized the \textbf{Third Generation} of computers and explain how it changed computer accessibility. \hfill (4 marks)
        \item List and explain \textbf{any three} computer generations (excluding the Third Generation which was discussed in part (a)). For each generation, state the key technology and one characteristic. \hfill (6 marks)
    \end{enumerate}

    \item \textbf{System Development Life Cycle (SDLC)} \hfill (10 marks)
    \begin{enumerate}[label=(\alph*)]
        \item Define the System Development Life Cycle (SDLC). Why is this important in software development? \hfill (2 marks)
        \item List and explain at least four phases of the SDLC Life Cycle. \hfill (8 marks)
    \end{enumerate}

    \item \textbf{Computer Literacy and Professionals} \hfill (10 marks)
    \begin{enumerate}[label=(\alph*)]
        \item Differentiate between \texttt{computer literacy} and \texttt{computer competency}. \hfill (4 marks)
        \item Differentiate between a \texttt{computer professional} and a \texttt{computer user}. Give two examples of each. \hfill (6 marks)
    \end{enumerate}

    \item \textbf{Number Systems and Storage} \hfill (10 marks)
    \begin{enumerate}[label=(\alph*)]
        \item Convert the binary number \textbf{1101101$_2$} to decimal. \hfill (3 marks)
        \item Convert the octal number \textbf{647$_8$} to binary. \hfill (3 marks)
        \item Convert \textbf{3.5 GB} to MB. Show your working. \hfill (4 marks)
    \end{enumerate}
\end{enumerate}

\newpage

\section*{SECTION B (50 marks)}
\textit{Answer ALL questions in this section.}

\begin{enumerate}[left=0pt, label=\bfseries Q\arabic*.]
    \item \textbf{Algorithms and Flowcharts} \hfill (15 marks)
    \begin{enumerate}[label=(\alph*)]
        \item Define an algorithm and list atleast four characteristics of a good algorithm. \hfill (5 marks)
        \item Write an algorithm to calculate the \textbf{volume} and \textbf{total surface area} of a cylinder given the radius and height. \hfill (5 marks)
        \emph{Formulas:}
        \begin{itemize}
            \item Volume $= \pi r^2 h$
            \item Total Surface Area $= 2\pi r^2 + 2\pi r h$
        \end{itemize}
        (Use $\pi = 3.142$)
        \item Draw a \textbf{flowchart} to represent the algorithm for determining whether a student qualifies for university admission based on their credit score. Use the following criteria: \hfill (5 marks)
        \begin{itemize}
            \item If credits $\ge 5$, print ``Qualify to be considered for admission''
            \item Else, print ``Not qualify to be considered for admission''
        \end{itemize}
    \end{enumerate}

    \item \textbf{Python Programming I} \hfill (15 marks)
    \begin{enumerate}[label=(\alph*)]
        \item Write a Python program that calculates the \textbf{area and perimeter} of a rectangle. The program should: \hfill (7 marks)
        \begin{enumerate}
            \item Ask the user for the \texttt{length} and \texttt{width}.
            \item Calculate and display the \texttt{area} ($\text{length} \times \text{width}$).
            \item Calculate and display the \texttt{perimeter} ($2 \times (\text{length} + \text{width})$).
            \item Use meaningful variable names and formatted strings for output.
        \end{enumerate}
        \item Write a Python program that takes a sentence from the user and performs the following string operations: \hfill (8 marks)
        \begin{enumerate}
            \item Displays the sentence in \textbf{uppercase}.
            \item Displays the sentence in \textbf{lowercase}.
            \item Displays the \textbf{length} of the sentence.
            \item Displays the \textbf{first 5 characters} of the sentence.
            % \item Displays the \textbf{last 5 characters} of the sentence.
        \end{enumerate}
    \end{enumerate}

    \item \textbf{Systems Analysis and Design} \hfill (20 marks)
    \begin{enumerate}[label=(\alph*)]
        \item List the \textbf{six key elements} of a system. \hfill (6 marks)
        \item Explain \textbf{five reasons} why information systems fail. \hfill (5 marks)
        \item Discuss \textbf{four reasons} why systems should be documented. \hfill (4 marks)
        \item Differentiate between \textbf{system software} and \textbf{application software}. Give \textbf{two examples} of each. \hfill (5 marks)
    \end{enumerate}
\end{enumerate}

\newpage

\section*{SECTION C (50 marks)}
\textit{Answer ALL questions in this section.}

\begin{enumerate}[left=0pt, label=\bfseries Q\arabic*.]
    \item \textbf{8-bit Binary Arithmetic} \hfill (15 marks) \\
    Assume an 8-bit binary machine where the \texttt{8th bit} is the sign bit (0 for positive, 1 for negative). Only 7 bits are available for magnitude.
    Given:
    \begin{itemize}
        \item $\mathbf{A = 25_{10}}$
        \item $\mathbf{B = 2F_{16}}$
        \item $\mathbf{C = 43_{8}}$
        \item $\mathbf{D = -14_{10}}$
    \end{itemize}
    Answer the following:
    \begin{enumerate}[label=(\alph*)]
        \item Express \texttt{D} as an 8-bit number using \texttt{2's complement} representation. Show your working. \hfill (3 marks)
        \item Calculate the binary sum of A and B. Provide the result in 8-bit binary. \hfill (4 marks)
        \item Calculate \textbf{A - B} using the \texttt{2's complement} method. Provide the result in 8-bit binary. \hfill (4 marks)
        \item Can the number \textbf{-150$_{10}$} be represented in this 8-bit machine? Justify your answer.\hfill (4 marks)
    \end{enumerate}

    \item \textbf{Python Programming II} \hfill (15 marks)
    \begin{enumerate}[label=(\alph*)]
        \item Write a Python program that acts as a \texttt{temperature converter}. The program should: \hfill (8 marks)
        \begin{enumerate}
            \item Ask the user for a temperature value.
            \item Ask the user whether to convert from \texttt{Celsius to Fahrenheit} or \texttt{Fahrenheit to Celsius}.
            \item Perform the appropriate conversion using the formulas:
            \begin{itemize}
                \item $^\circ\mathrm{F} = (^\circ\mathrm{C} \times 9/5) + 32$
                \item $^\circ\mathrm{C} = (^\circ\mathrm{F} - 32) \times 5/9$
            \end{itemize}
            \item Display the result with an appropriate message.
        \end{enumerate}
        \item Write a Python program that uses a \texttt{loop} to generate the multiplication table for a number provided by the user, from \texttt{1 to 10 times}. \hfill (7 marks) \\
        \emph{Example output (if user enters 7):}
        \begin{verbatim}
        7 × 1 = 7
        7 × 2 = 14
        ...
        7 × 10 = 70
        \end{verbatim}
    \end{enumerate}

    \item \textbf{Computer Architecture and Number Systems} \hfill (20 marks)
    \begin{enumerate}[label=(\alph*)]
        \item Explain the \texttt{stored program concept} in the Von Neumann architecture and state its significance. \hfill (5 marks)
        \item Differentiate between \texttt{RAM} and \texttt{ROM} in terms of: \hfill (5 marks)
        \begin{itemize}
            \item Volatility
            \item Usage
            \item Read/Write capabilities
        \end{itemize}
        \item Describe the role of the Control Unit (CU) in a CPU and explain how it coordinates with other components. \hfill (5 marks)
        \item Convert the following: \hfill (5 marks)
        \begin{enumerate}
            \item The hexadecimal number \textbf{3A7$_{16}$} to binary. \hfill (3 marks)
            \item The decimal number \textbf{158$_{10}$} to octal. \hfill (2 marks)
        \end{enumerate}
    \end{enumerate}
\end{enumerate}

\vspace{1cm}
\begin{center}
    \textbf{END OF EXAMINATION}
\end{center}

\end{document}